\textsc{[Greedy]}
\\
You're working on a
long research paper, and your draconian library will only allow you to have
$k$ books checked out at once. 
You know that you’ll need more
than this over the course of working on the paper, but at any point in time,
you'd like to have ready access to the $k$ 
books that are most relevant at
that time. How should you decide which books to check out, and when should
you return some in exchange for others, to minimize the number of times you
have to exchange a book at the library?

As a concrete example, 
suppose $k = 2$ (i.e., you can only 
check out two books at a time.
Suppose you need to refer to 
three books $a, b, c$.
Furthermore, suppose you need to 
use these books in this order:
\[
a,b,c,b,c,a,b
\]
and suppose that you have books $a$ on day 0
and $b$ on day 1.
Then on the day 2, since you need book $c$,
you need to return either book $a$ or book $b$
in order to check out book $c$.
Suppose you return book $a$.
Therefore on day 2, you have books $b$ and $c$.
On day 3, you need book $b$, which you do have.
On day 4, you need book $c$, which is
already checked out.
So on day 3 and 4, you do not need to return
any books.
But on day 5, you need book $a$.
This means that you need to return either
book $b$ or $c$.
Etc.
The goal is (obviously) to prevent trips to the library.

Design an algorithm that minimizes library trips.
Implement your algorithm in C\texttt{++}.
Here's a sample execution:
\begin{Verbatim}[frame=single,commandchars=\\\{\},fontsize=\footnotesize]
\userinput{2 3 7 0 1 2 1 2 0 1}
1. 0->2
2.
3.
4. 2->0
5.
\end{Verbatim}
The first number, i.e., $2$, in the input is the number of books you can check out.
The second number, i.s., $3$, is the number of books you need to use.
The third number, i.e., $7$, is the number of days you will be doing research.
The rest, \verb!0 1 2 1 2 0 1!, is of (7) accesses to the books.
Note that the books are numbered 0, 1, 2 if there are 3 books.
The rest describes visits to the library.
At the beginning, you already have books 0 and 1.
So the output \verb!1. 0->2! says that you return book 0 to check out 2.
\verb!2.! says that you do not visit the library.
This is followed by \verb!3.! which again says that you do not visit the library.
\verb!4. 2->1! says that you visit the library to return book 2 in order to check out book 1.
Etc.


[\textsc{Aside}. This problem occurs in every situation
where you need to work with more data than what a resource can handle
at one time -- i.e., there's a memory hierarchy.
In the above library example, you can only hold two books at one time.
The library can hold more.
In the case of computer architecture, main memory and hard drive.
It's better to have data in main memory
because that's faster than hard drive. But the problem is that
your computer has a lot less capacity in main memory than the hard drive.
For modern computer architecture, there's on-chip cache memory which is
faster than main memory.
Likewise, a database system stores pages of data in main memory
to create a buffer pool and attempts to predict what pages of data stored
in the hard drive in order to prevent replacing buffer pool pages with
new pages from the hard drive.
In the above library example, when you return a book, say $a$, already checked out
for a new book, say $c$, we say that your algorithm is \defterm{evicting} $a$.
Sometimes $a$ is also called a \defterm{victim}.
The overall algorithm that exchanging a piece of data in cache with another
is frequently called a \defterm{replacement algorithm}.
Note that for our library problem, it assumes that you know ahead of time
what book will be needed in the future. In real life, frequently
only an approximate access pattern is known, i.e., the access
pattern is given in terms of probabilistic measures.]
